\documentclass[aspectratio=169,11pt]{beamer}
% Add `handout' to the class options to collapse overlays into one page per frame.
\usetheme[accent=gray]{Gondolin}
% Options: \usetheme[sans]{Gondolin}          Alegreya Sans body text
%          \usetheme[accent=aqua]{Gondolin}   gruvbox accent (indigo, red, orange, yellow,
%                                            green, aqua, blue, purple, gray)
%          \usetheme[nodividers]{Gondolin}    no dark divider slide per section
%          \usetheme[sectionpages]{Gondolin}  outline slide at each section
%          \usetheme[noframeno]{Gondolin}     no frame numbers
%          \usetheme[sharednumbering]{Gondolin}  Fact 1, Remark 2, ... (one counter)
%          \usetheme[nonatbib]{Gondolin}      bring your own biblatex
%
% Speaker notes: uncomment one of these.
% \setbeameroption{show notes on second screen=right}  % for pdfpc
% \setbeameroption{show only notes}                    % printable script

\title{An Example Talk with the Gondolin Theme}
\subtitle{Montpellier's layout, cosmic latte and gruvbox, Arguelles' fonts}
\author{First Author}
\institute{International Institute of Information Technology, Hyderabad}
\email{first.author@example.com}
\event{Course Seminar}
\date{\today}

\begin{document}

\titleframe

\begin{frame}{Outline}
  \tableofcontents
\end{frame}

\section{Introduction}

\subsection{Motivation}

\begin{frame}{Lists and overlays}{a frame subtitle}
  \begin{itemize}
    \item<1-> The headline shows the section, the subsection and the talk title.
    \item<2-> Overlays work as usual; a handout build collapses them.
      \begin{itemize}
        \item Second-level items are dimmer.
      \end{itemize}
    \item<3-> \alert{Alerted text} is gruvbox red.
  \end{itemize}
  \vspace{1em}
  \begin{enumerate}
    \item Citations are author--year: \citet{lamport94}, \citep{knuth84}.
    \item Links take the accent colour: \url{https://ctan.org}.
  \end{enumerate}
  \note{Speaker notes go in \texttt{\textbackslash note}. They appear only when
    notes are switched on in the preamble.}
\end{frame}

\subsection{Mathematics}

\begin{frame}{Maths}
  Let $\mathbf{X} \in \mathbb{R}^{n \times d}$ be a design matrix and
  $\mathbf{y} \in \mathbb{R}^{n}$ a target vector. For $\lambda > 0$ the ridge
  estimate is
  \[
    \hat{\mathbf{w}} = \operatorname*{arg\,min}_{\mathbf{w} \in \mathbb{R}^{d}}
      \lVert \mathbf{y} - \mathbf{X}\mathbf{w} \rVert_2^2
      + \lambda \lVert \mathbf{w} \rVert_2^2
    = (\mathbf{X}^\top \mathbf{X} + \lambda \mathbf{I})^{-1} \mathbf{X}^\top \mathbf{y},
  \]
  where $\mathbf{I}$ is the $d \times d$ identity matrix.
\end{frame}

\begin{frame}{Theorems}
  \begin{definition}
    A matrix $\mathbf{A} \in \mathbb{R}^{d \times d}$ is \emph{positive definite} if
    $\mathbf{v}^\top \mathbf{A} \mathbf{v} > 0$ for every nonzero
    $\mathbf{v} \in \mathbb{R}^{d}$.
  \end{definition}
  \begin{theorem}
    For every $\lambda > 0$ the matrix $\mathbf{X}^\top \mathbf{X} + \lambda \mathbf{I}$
    is positive definite, hence invertible.
  \end{theorem}
  \begin{remark}
    Each kind of block is numbered on its own: this is Remark 1.
  \end{remark}
  \begin{proof}
    For nonzero $\mathbf{v}$ we have
    $\mathbf{v}^\top (\mathbf{X}^\top \mathbf{X} + \lambda \mathbf{I}) \mathbf{v}
     = \lVert \mathbf{X}\mathbf{v} \rVert_2^2 + \lambda \lVert \mathbf{v} \rVert_2^2
     \ge \lambda \lVert \mathbf{v} \rVert_2^2 > 0$.
  \end{proof}
\end{frame}

\section{Results}

\begin{frame}{Blocks}
  \begin{columns}[T]
    \column{.32\textwidth}
      \begin{block}{Block}
        A plain block.
      \end{block}
    \column{.32\textwidth}
      \begin{alertblock}{Alert block}
        Something to watch out for.
      \end{alertblock}
    \column{.32\textwidth}
      \begin{exampleblock}{Example block}
        An example; its title takes the accent colour.
      \end{exampleblock}
  \end{columns}
\end{frame}

\begin{frame}[fragile]{Code}
  Code frames must be \texttt{[fragile]}. Inline code: \lstinline|np.linalg.solve|.
\begin{lstlisting}[language=Python]
import numpy as np

def ridge(X, y, lam):
    """Closed-form ridge estimate."""
    d = X.shape[1]
    return np.linalg.solve(X.T @ X + lam * np.eye(d), X.T @ y)  # shape (d,)
\end{lstlisting}
\end{frame}

\begin{frame}{A figure}
  \begin{figure}
    \centering
    \includegraphics[height=.5\textheight]{example-image}
    \caption{Captions are numbered.}
  \end{figure}
\end{frame}

\begin{frame}{Image and text}
  \imagetext[caption={The image goes on the right by default.}]{example-image}{%
    \begin{itemize}
      \item \texttt{\textbackslash imagetext\{file\}\{text\}} puts the text on the left.
      \item Options: \texttt{left}, \texttt{width=.45}, \texttt{height=.75},
            \texttt{caption=...}, \texttt{align=T}.
    \end{itemize}}
\end{frame}

\begin{frame}{Image on the left}
  \imagetext[left, width=.35]{example-image-b}{%
    With \texttt{left} the image moves to the left; \texttt{width=.35} makes its
    column narrower, and the text column takes the rest.}
\end{frame}

\begin{frame}[bg=example-image-a]{An image frame}
  \texttt{[bg=file]} fills the slide with an image, cropped to fit.
\end{frame}

\begin{frame}[standout]
  A \emph{standout} frame for the one sentence\\ people should remember.
\end{frame}

\section{Conclusion}

\begin{frame}[fragile]{Summary}
  \begin{itemize}
    \item Drop \texttt{beamerthemeGondolin.sty} next to your slides.
    \item Use \verb|\documentclass[aspectratio=169]{beamer}| and \verb|\usetheme{Gondolin}|.
  \end{itemize}
\end{frame}

\begin{frame}[standout]
  Thank you!\\[.5em]
  {\normalsize Questions?}
\end{frame}

% Everything after \appendix is a backup slide: it gets a divider and does not
% count towards the "n / N" total.
\appendix

\nocite{goodfellow16}
\begin{frame}{References}
  \bibliographystyle{rivendell}
  \bibliography{refs}
\end{frame}

\begin{frame}{A backup slide}
  Details you only show if someone asks.
\end{frame}

\end{document}
